\documentclass[11pt]{amsart}
\usepackage[T1]{fontenc}
\usepackage[utf8]{inputenc}
\usepackage[british]{babel}
\usepackage[margin=1in]{geometry}
\parindent=8pt
\parskip=3pt
\linespread{1.1}
\usepackage{amsmath,amssymb,amsthm}
\usepackage{palatino,euler,microtype}
% Euler expects the upright shape name used by the reference document.
\renewcommand{\updefault}{n}
\usepackage{booktabs,longtable,array}
\usepackage{graphicx}
\usepackage{tikz}
\usetikzlibrary{positioning,fit,backgrounds,calc,shapes.geometric}
\input{figures/jfef-grafos.tikzstyles}
\usepackage{xcolor}
\usepackage{float,placeins}
\usepackage{needspace}
\definecolor{darkblue}{rgb}{0,0,0.7}
\definecolor{darkred}{rgb}{0.7,0,0}
\newcommand{\defi}[1]{\textit{\color{blue}#1}}
\usepackage[pagebackref=true]{hyperref}
\hypersetup{colorlinks=true,linkcolor=darkred,citecolor=darkblue,urlcolor=darkblue,
pdftitle={Counterexamples and enumeration for contractible and dismantlable graphs}}
\setlength{\emergencystretch}{2em}



\usepackage{enumitem,siunitx}
\sisetup{group-separator={,},group-minimum-digits=4}
\newcommand{\I}{\mathfrak I}
\newcommand{\IS}{\mathfrak I_S}
\newcommand{\ISV}{\mathfrak I_{SV}}
\newcommand{\Dis}[1]{#1\text{-Dis}}
\newcommand{\ZZ}{\mathbb Z}
\newcommand{\Ztwo}{\mathbb F_2}
\newcommand{\cl}{\operatorname{cl}}
\newcommand{\Nb}{N}
\newcommand{\Gt}{G_{13}}
\newcommand{\Ga}{G_{14}^{a}}
\newcommand{\Gb}{G_{14}^{b}}
\newcommand{\Gc}{G_{14}^{c}}
\newcommand{\Dunce}{D_{15}}
\newtheorem{theorem}{Theorem}[section]
\newtheorem{lemma}[theorem]{Lemma}
\newtheorem{proposition}[theorem]{Proposition}
\newtheorem{corollary}[theorem]{Corollary}
\newtheorem{definition}[theorem]{Definition}
\newtheorem{remark}[theorem]{Remark}
\newtheorem{construction}[theorem]{Construction}
\newtheorem{example}[theorem]{Example}
\newtheorem{observation}[theorem]{Observation}
\newtheorem{conjecture}[theorem]{Conjecture}
% Numerical values used in this manuscript. All row-14 values are K4-free.
\newcommand{\RowThirteenAcyclic}{\num{5619022190286}}
\newcommand{\RowThirteenI}{\num{5619022190283}}
\newcommand{\RowThirteenIS}{\num{5619022190283}}
\newcommand{\RowFourteenAcyclic}{\num{755877673752}}
\newcommand{\RowFourteenI}{\num{755877673127}}
\newcommand{\RowFourteenIS}{\num{755877673119}}
\newcommand{\RowFourteenCone}{\num{20797002}}
\newcommand{\NumTwoDisNotOneDisThirteen}{\num{35757}}

\renewcommand{\urladdrname}{\textit{ORCID}}
\title[Counterexamples and enumeration for contractible graphs]{Counterexamples and enumeration for contractible and dismantlable graphs}
\author[J. F. Espinoza]{Jes\'us F. Espinoza\textsuperscript{*}}
\thanks{\textsuperscript{*}Corresponding author}
\author[M. E. Fr\'ias-Armenta]{Mart\'in Eduardo Fr\'ias-Armenta}
\author[R. G. Hern\'andez-Amador]{Rosal\'ia G. Hern\'andez-Amador}
\author[H. A. Hern\'andez-Hern\'andez]{H\'ector A. Hern\'andez-Hern\'andez}
\address{Departamento de Matem\'aticas, Universidad de Sonora, Hermosillo, Sonora, M\'exico}
\email[Jes\'us F. Espinoza]{jesusfrancisco.espinoza@unison.mx}
\urladdr[Jes\'us F. Espinoza]{https://orcid.org/0000-0001-5601-3634}
\email[Mart\'in Eduardo Fr\'ias-Armenta]{eduardo.frias@unison.mx}
\urladdr[Mart\'in Eduardo Fr\'ias-Armenta]{https://orcid.org/0000-0002-7687-1270}
\email[Rosal\'ia G. Hern\'andez-Amador]{rosalia.hernandez@unison.mx}
\urladdr[Rosal\'ia G. Hern\'andez-Amador]{https://orcid.org/0000-0001-9657-7012}
\email[H\'ector A. Hern\'andez-Hern\'andez]{hector.hernandez@unison.mx}
\urladdr[H\'ector A. Hern\'andez-Hern\'andez]{https://orcid.org/0000-0002-4046-4841}
\date{}
\subjclass[2020]{Primary 05C75; Secondary 05C30, 05C85, 05E45, 55U10, 57Q10}
\keywords{contractible transformations, strong I-contractibility, dismantlability,
clique complex, odd torsion, enumeration}
\begin{document}
\begin{abstract}
Contractible transformations of graphs preserve the homotopy type of their
clique complexes. Restricting to deletions, and imposing this restriction recursively
on their neighbourhoods, gives strong I-contractibility. We construct
$13$-vertex graphs that are acyclic over $\mathbb F_2$ but have first
integral homology $\mathbb Z/3$, and $14$-vertex I-contractible graphs
outside the strong class. Adding one admissible edge to each of the latter
graphs produces a strong graph. These examples refute Conjectures 14 and 15
of Dochtermann, Espinoza, Fr\'ias-Armenta and Hern\'andez.
We give exact tests for acyclicity, strong contractibility and
dismantlability, prove their correctness, and extend the cardinality
table for connected graphs through order $13$ and connected
$K_4$-free graphs of order $14$.
The table distinguishes integral from mod-$2$ acyclicity and establishes
minimum orders $13$ and $14$ for the two separations. There are exactly
eight $K_4$-free graphs of order $14$ in the non-strong I-contractible
class; their minimum number of edges is $44$.
\end{abstract}
\maketitle
\section{Introduction}\label{sec:intro}

A contractible graph can be transformed into a single vertex by adding
and deleting vertices and edges whose corresponding neighbourhoods are
contractible. Requiring a reduction by deletions alone, and imposing the
same requirement recursively on the neighbourhoods, gives the class of
strong I-contractible graphs. What is the least order of an I-contractible graph that is not strong
I-contractible? Does replacing contractibility by acyclicity over
$\mathbb F_2$ lower this minimum?

Ivashchenko introduced contractible transformations of graphs
\cite{Ivashchenko1993,Ivashchenko1994a}. The clique complex $\cl(G)$ has
the cliques of $G$ as its simplices. The class $\I$ of I-contractible
graphs is precisely the class of graphs with contractible clique complex
\cite[Corollary~1.5]{DochtermannMatsushita2025}.
Strong I-contractible graphs, denoted by $\IS$, have collapsible clique
complexes \cite[Theorem~1]{EFH2022}; hence $\IS\subseteq\I$.

Dochtermann, Espinoza, Fr\'ias-Armenta and Hern\'andez
\cite{DEFH2023} counted several contractible and dismantlable families
through order $11$. They also exhibited a $15$-vertex graph $\Dunce$ whose
clique complex is a flag triangulation of the dunce hat, with
$\Dunce\in\I\setminus\IS$, and a one-edge extension $\Dunce^+\in\IS$.
In our notation, their conjectures assert that these examples have
minimum order:

\begin{quote}
\textbf{Conjecture 14 of \cite[p.~8]{DEFH2023}.}
The graph $\Dunce$ is smallest by number of vertices among graphs with
trivial reduced $\mathbb F_2$-homology that are not strong
I-contractible. In particular, it is smallest among the graphs in
$\I\setminus\IS$.

\smallskip
\textbf{Conjecture 15 of \cite[p.~9]{DEFH2023}.}
The graph $\Dunce^+$ is smallest by number of vertices among graphs
$G\in\IS$ having an edge $uv$ for which the common neighbourhood
$N_G(u,v)$ belongs to $\IS$ but $G-uv\notin\IS$.
\end{quote}

We refute both conjectures. Three graphs on $13$ vertices have clique
complexes homotopy equivalent to the Moore space $M(\ZZ/3,1)$.
Their reduced homology vanishes over $\mathbb F_2$, whereas their first
integral homology is $\ZZ/3$. They therefore lie outside $\I$ and
refute the first assertion of Conjecture~14. A graph on $14$ vertices
belongs to $\I\setminus\IS$ and refutes its second assertion. Its
membership in $\I$ follows from one admissible edge addition followed
by a strong reduction. The same extension refutes Conjecture~15.
Section~\ref{sec:examples} gives the constructions and their proofs.

Our second contribution is an extension of the cardinality table,
together with algorithms and proofs that justify its computation.
We count connected graphs up to isomorphism in the contractible and
dismantlable families through order $13$, and in the connected
$K_4$-free universe at order $14$. Integral and $\mathbb F_2$
acyclicity are recorded separately. The integral and I-contractible
columns are derived from the $\mathbb F_2$ census and the classification
of its non-strong graphs. The three torsion examples are
exactly the $\mathbb F_2$-acyclic graphs outside $\IS$ through order
$13$. Consequently, the minimum order for the first assertion of
Conjecture~14 is $13$, while the minimum for its second assertion and
for the property in Conjecture~15 is $14$. At order $14$ there are
exactly eight $K_4$-free members of $\I\setminus\IS$; the fewest
edges among these eight is $44$.

The table also records the hierarchy of $k$-dismantlable graphs
introduced by Fieux and Jouve \cite{FieuxJouve2020}.
Throughout the listed universes, $2$-Dis and $3$-Dis coincide with
the strong vertex I-contractible class $\ISV$. The first separation
between $1$-Dis and $2$-Dis occurs at order $12$, with $20$ graphs.

The computation uses exact recursions for the strong classes and the
dismantling levels. At each stage these recursions consider all possible
first deletions. The remaining acyclic graphs are classified by explicit
contractible transformations and integral homology. For $K_4$-free
graphs, vertex neighbourhoods have clique complexes of dimension at most one and
edge neighbourhoods are discrete; this gives local admissibility criteria
and reduces the order-$14$ classification to graphs with no
admissible deletion.

The $13$-vertex construction also attains the minimum order for odd
torsion in the first integral homology of a flag complex, by comparison
with Adamaszek's classification \cite{Adamaszek2014}.
The three $14$-vertex graphs constructed below have contractible,
non-collapsible clique complexes. The computer-assisted classification through order
$13$ shows that $14$ is the minimum number of vertices of such a
flag complex (Corollary~\ref{cor:minimality}). Earlier examples arise
from Bing's house \cite{ChenYauYeh2001} and from flag triangulations
of the dunce hat \cite{FieuxJouve2020,DEFH2023}.
Our classification of order-$14$
graphs is restricted to the $K_4$-free class; the minimum by number
of vertices follows from the unrestricted census through order $13$.

\medskip\noindent\textbf{Organization.}
Section~\ref{sec:prelim} fixes the graph operations and homological
conventions. Section~\ref{sec:examples} proves the refutations using
explicit examples. Sections~\ref{sec:algorithms} and~\ref{sec:searches}
justify the enumeration and derive the table and the minimum orders.
The appendix supplies the finite mathematical data used in the
constructions.

\medskip\noindent\textbf{Conventions and notation.}
Graphs are finite and simple. Their order is their number of vertices.
Homology refers to the clique complex and is reduced unless stated
otherwise. The empty graph belongs to none of the contractible or
acyclic classes considered here. Cardinalities count isomorphism classes.

\section{Graph reductions and homology}\label{sec:prelim}

The graph operations below are determined by induced neighbourhoods.
For a vertex $v$ of $G$, we write $N_G(v)$ for the subgraph induced by
its neighbours. For distinct vertices $u,v$, the common neighbourhood
$N_G(u,v)$ is induced by the vertices adjacent to both. We write $G-v$
for a vertex deletion, $G-uv$ for an edge deletion, and $G+uv$ for the
addition of a missing edge. The clique complex $\cl(G)$ consists of the
cliques of $G$; a complex of this form is called \defi{flag}.
The clique complexes of $N_G(v)$ and $N_G(u,v)$ are the links of $v$
and, when $uv$ is an edge, of $uv$ in $\cl(G)$.

\subsection{Contractible and strong contractible graphs}
\label{sec:contractible}
The two classes differ in the permitted operations and in the
condition imposed on neighbourhoods.

\begin{definition}[{\cite[Definition~1]{DEFH2023}}]\label{def:I}
The class $\I$ of \defi{I-contractible graphs} is the least class
containing $K_1$ and closed under the following operations:
\begin{enumerate}[label=(\roman*),leftmargin=2.5em]
\item delete a vertex $v$ when $N_G(v)\in\I$;
\item add a new vertex whose neighbourhood is an induced subgraph
of $G$ belonging to $\I$;
\item delete an edge $uv$ when $N_G(u,v)\in\I$;
\item add a missing edge $uv$ when $N_G(u,v)\in\I$.
\end{enumerate}
The empty graph does not belong to $\I$.
\end{definition}

The operations in Definition~\ref{def:I} preserve the simple homotopy
type of the clique complex \cite[Theorem~2.8]{DochtermannMatsushita2025};
the graph formulation goes back to Ivashchenko and to Chen, Yau and Yeh
\cite{Ivashchenko1994a,ChenYauYeh2001}. In particular, every member
of $\I$ has contractible clique complex. The converse holds as well
\cite[Corollary~1.5]{DochtermannMatsushita2025}.

\begin{definition}[{\cite[Definition~3 and Lemma~1]{EFH2022};
\cite[Definitions~4--5]{DEFH2023}}]\label{def:IS}
The class $\IS$ of \defi{strong I-contractible graphs} contains $K_1$
and is defined recursively, for graphs of order at least two, by
\[
 G\in\IS\quad\Longleftrightarrow\quad
 \begin{cases}
 G-v,\ N_G(v)\in\IS&\text{for some }v\in V(G),\quad\text{or}\\
 G-uv,\ N_G(u,v)\in\IS&\text{for some }uv\in E(G).
 \end{cases}
\]
The class $\ISV$ of \defi{strong vertex I-contractible graphs} has the
same base case and uses only the first alternative, with both smaller
graphs required to belong to $\ISV$. Neither class contains the empty
graph.
\end{definition}

Each recursive argument has smaller $|V|+|E|$, so these definitions are
well founded. Reversing a successful reduction gives the equivalent
definitions by gluings in the cited papers. Induction on $|V|+|E|$
gives
\begin{equation}\label{eq:strong-inclusions}
 \ISV\subseteq\IS\subseteq\I.
\end{equation}
For the second inclusion, the remaining graph and the tested
neighbourhood lie in $\I$ by induction, and the inverse operation
reconstructs the original graph by an allowed gluing.

An operation is \defi{$\IS$-admissible} if its tested neighbourhood
belongs to $\IS$. Unless a different class is named, admissibility
will mean $\IS$-admissibility. A \defi{strong reduction} is a sequence
of admissible deletions ending at $K_1$. The empty-neighbourhood
convention means that neither an isolated vertex nor an edge with empty
common neighbourhood can be deleted. Every nonempty tree belongs to
$\ISV$: deleting a leaf tests the neighbourhood $K_1$ and reduces the
order. Conversely, a triangle-free member of $\I$ is a tree, since
its clique complex is connected, contractible and of dimension at most one.

Membership in $\IS$ requires the existence of a successful reduction.
A different admissible first deletion can destroy that possibility,
as the examples for Conjecture~15 will show. Thus the recursive test
must consider all possible first deletions.

\subsection{Dismantling levels}
Vertex reductions also give a hierarchy in which the condition on the
neighbourhood changes with the level.

\begin{definition}[{\cite[Definition~1]{FieuxJouve2020}}]\label{def:kdis}
A \defi{cone} is a nonempty graph with a vertex adjacent to every other
vertex. Cones form the class $(-1)$-Dis. For $k\geq0$, a vertex $v$
is \defi{$k$-dismantlable} when $N_G(v)$ belongs to $(k-1)$-Dis.
A graph is \defi{$k$-dismantlable}, or belongs to $k$-Dis, if it can
be reduced to $K_1$ by successive deletions of $k$-dismantlable
vertices. The singleton belongs to every level, and the empty graph
to none.
\end{definition}

At level zero, the condition says that some neighbour $u$ of $v$ is
adjacent to all the other neighbours of $v$. Equivalently,
$N_G[v]\subseteq N_G[u]$, where brackets denote closed neighbourhoods.
The hierarchy satisfies
\begin{equation}\label{eq:dismantling-inclusions}
 (-1)\text{-Dis}\subseteq0\text{-Dis}\subseteq1\text{-Dis}
 \subseteq\cdots\subseteq\ISV\subseteq\IS\subseteq\I;
\end{equation}
see \cite[Proposition~2]{FieuxJouve2020} and
\cite[Lemma~19 and Theorem~23]{DEFH2023}.
The union of the finite dismantling levels is $\ISV$
\cite[Theorem~23]{DEFH2023}.

\subsection{Coefficient choices}
We use homology to distinguish graphs whose clique complexes are
contractible from graphs with the same homology over $\mathbb F_2$.
Write $\widetilde H_i(G;R)=\widetilde H_i(\cl(G);R)$.
For $R=\ZZ$ or $\mathbb F_2$, let $\mathcal A_R$ be the class of
nonempty graphs for which $\widetilde H_i(G;R)=0$ for every $i\geq0$.
The Euler characteristic is
\[
 \chi(G)=\sum_{i\geq0}(-1)^if_i(G),
\]
where $f_i(G)$ counts the cliques with $i+1$ vertices.

For a nonempty finite complex $K$, the universal coefficient theorem
\cite[Theorem~3A.3]{Hatcher2002} gives, for $i\geq0$,
\begin{equation}\label{eq:uct}
 0\longrightarrow\widetilde H_i(K;\ZZ)\otimes\mathbb F_2
 \longrightarrow\widetilde H_i(K;\mathbb F_2)
 \longrightarrow\operatorname{Tor}^{\ZZ}_1
       (\widetilde H_{i-1}(K;\ZZ),\mathbb F_2)
 \longrightarrow0.
\end{equation}
We use the augmented convention $\widetilde H_{-1}(K;\ZZ)=0$.
Consequently,
\begin{equation}\label{eq:homology-inclusions}
 \IS\subseteq\I\subseteq\mathcal A_{\ZZ}
 \subseteq\mathcal A_{\mathbb F_2}.
\end{equation}
For a connected finite complex, $\mathbb F_2$-acyclicity allows finite
integral homology of odd order. In particular, both
$(\ZZ/3)\otimes\mathbb F_2$ and
$\operatorname{Tor}^{\ZZ}_1(\ZZ/3,\mathbb F_2)$ vanish.
This accounts for the different hypotheses in the two assertions of
Conjecture~14. Throughout the article, unqualified acyclicity means
membership in $\mathcal A_{\mathbb F_2}$.

\subsection{A disc with boundary identifications}\label{sec:discs}
The first construction is most easily understood by obtaining its
clique complex from a triangulated disc. The following criterion
determines the homotopy type from the boundary identifications.

\begin{proposition}\label{prop:disc}
Let $D$ be a triangulated closed disc and $q:D\to X$ a simplicial
map. Suppose that $q$ is injective on the interior of $D$ and identifies
$X$ with the quotient of $D$ obtained solely by identifying boundary
points having the same image. If $S=q(\partial D)$ is a connected
graph with exactly one cycle, then
\[
 X\simeq S^1\cup_d D^2,
\]
where $d$ is the degree of the boundary map after contracting a maximal
tree in $S$. For $|d|=1$ the complex is contractible. For $|d|=3$
it is homotopy equivalent to $M(\ZZ/3,1)$ and has
$\widetilde H_1(X;\ZZ)\cong\ZZ/3$ and zero reduced homology in all
other degrees.
\end{proposition}
\begin{proof}
The hypotheses identify $X$ with the adjunction space $S\cup_fD^2$.
Collapsing a maximal tree $T$ in $S$ is a homotopy equivalence, since
$T$ is a contractible subcomplex \cite[Propositions~0.16--0.17]{Hatcher2002}.
The quotient is $(S/T)\cup_{pf}D^2$, with $S/T\cong S^1$.
Maps $S^1\to S^1$ are classified up to homotopy by their degree
\cite[Theorem~1.7]{Hatcher2002}, and homotopic attaching maps have
homotopy equivalent adjunction spaces
\cite[Proposition~0.18]{Hatcher2002}. This proves the asserted
homotopy type. For $|d|=1$ the map is homotopic to a homeomorphism,
whose mapping cone is contractible. For $|d|=3$ this is the defining
adjunction space of $M(\ZZ/3,1)$, and its cellular chain complex in
positive degrees is $0\to\ZZ\xrightarrow{d}\ZZ\to0$.
\end{proof}

\section{Counterexamples on thirteen and fourteen vertices}\label{sec:examples}

The two assertions in Conjecture~14 are refuted at different orders.
Odd integral torsion supplies a graph with trivial homology over
$\Ztwo$ outside $\IS$. A second construction belongs to $\I$ but admits
no first strong deletion. Adding one edge to that construction also
gives a counterexample to Conjecture~15.

Throughout this section, $\Gt$ has vertex set $\{0,\ldots,12\}$ and
$\Ga,\Gb,\Gc$ have vertex set $\{0,\ldots,13\}$. Their edges are defined
by Table~\ref{tab:witness-adjacency}. The figures use these same labels.

\subsection{Odd torsion on thirteen vertices}\label{sec:thirteen}

The clique complex of $\Gt$ is obtained from a triangulated disc by
wrapping its boundary three times around a four-cycle. This description accounts
for the torsion and gives a direct proof of the first refutation.

\begin{theorem}\label{thm:A}
Let $\Gt'=\Gt+\{4,7\}$ and $\Gt''=\Gt'+\{5,6\}$.
The graphs $\Gt,\Gt',\Gt''$ have $40$, $41$ and $42$ edges, respectively,
and their clique complexes are homotopy equivalent to $M(\ZZ/3,1)$.
In particular, they are $\Ztwo$-acyclic and none belongs to $\I$.
They refute the first assertion of Conjecture~14.
\end{theorem}

\begin{proof}
The graph $\Gt$ is $K_4$-free and has $28$ triangles. Each of its edges
lies in two triangles, except for the four edges of
$8\text{--}10\text{--}9\text{--}11\text{--}8$, which lie in three.
Cutting along that cycle gives the triangulated disc on the left of
Figure~\ref{fig:G13}. An explicit triangle list and an order of attachment
proving that it is a disc are given in
Appendix~\ref{app:disc-certificates}.
Its boundary labels, starting at the top and read clockwise, are
\[
 10,9,11,8\mid10,9,11,8\mid10,9,11,8.
\]
Each interior vertex occurs once. Identifying equally labelled boundary
vertices and their corresponding boundary edges recovers precisely
$\cl(\Gt)$.
Proposition~\ref{prop:disc}, with attaching degree $3$, therefore gives
\[
 \cl(\Gt)\simeq M(\ZZ/3,1),\qquad
 \widetilde H_1(\Gt;\ZZ)\cong\ZZ/3,
 \qquad \widetilde H_i(\Gt;\ZZ)=0\quad(i\ne1).
\]
The universal coefficient theorem gives vanishing reduced homology
over $\Ztwo$ in every degree.

For each of the two specified edge additions, the common neighbourhood
at the time of addition is the singleton $\{12\}$. Both additions are
admissible and preserve the homotopy type of the clique complex.
Nonzero integral homology excludes all three graphs from $\I$, hence
also from $\IS$. Since they have fewer than the $15$ vertices of the
graph in Conjecture~14, they refute its first assertion.
\end{proof}

\begin{figure}[tbp]
\centering
\begin{tikzpicture}[x=1cm,y=1cm]
\definecolor{witnessred}{RGB}{181,55,36}
\definecolor{witnessblue}{RGB}{32,89,140}
\coordinate (v0) at (0.942769700,-1.632925019);
\coordinate (v1) at (1.632925019,0.942769700);
\coordinate (v2) at (-1.632925019,-0.942769700);
\coordinate (v3) at (-0.942769700,1.632925019);
\coordinate (v4) at (1.981676324,-0.530988571);
\coordinate (v5) at (-0.530988571,-1.981676324);
\coordinate (v6) at (0.530988571,1.981676324);
\coordinate (v7) at (-1.981676324,0.530988571);
\coordinate (v8) at (-1.700000000,2.944486373);
\coordinate (v9) at (3.400000000,0.000000000);
\coordinate (v10) at (-1.700000000,-2.944486373);
\coordinate (v11) at (1.700000000,-2.944486373);
\coordinate (v12) at (1.700000000,2.944486373);
\coordinate (v13) at (-3.400000000,-0.000000000);
\coordinate (v14) at (-2.944486373,-1.700000000);
\coordinate (v15) at (2.944486373,-1.700000000);
\coordinate (v16) at (0.000000000,3.400000000);
\coordinate (v17) at (2.944486373,1.700000000);
\coordinate (v18) at (0.000000000,-3.400000000);
\coordinate (v19) at (-2.944486373,1.700000000);
\coordinate (v20) at (0.000000000,0.000000000);
\fill[blue!3] (v0) -- (v4) -- (v15) -- cycle;
\fill[blue!3] (v0) -- (v4) -- (v20) -- cycle;
\fill[blue!3] (v0) -- (v5) -- (v18) -- cycle;
\fill[blue!3] (v0) -- (v5) -- (v20) -- cycle;
\fill[blue!3] (v0) -- (v11) -- (v15) -- cycle;
\fill[blue!3] (v0) -- (v11) -- (v18) -- cycle;
\fill[blue!3] (v1) -- (v4) -- (v9) -- cycle;
\fill[blue!3] (v1) -- (v4) -- (v20) -- cycle;
\fill[blue!3] (v1) -- (v6) -- (v12) -- cycle;
\fill[blue!3] (v1) -- (v6) -- (v20) -- cycle;
\fill[blue!3] (v1) -- (v9) -- (v17) -- cycle;
\fill[blue!3] (v1) -- (v12) -- (v17) -- cycle;
\fill[blue!3] (v2) -- (v5) -- (v10) -- cycle;
\fill[blue!3] (v2) -- (v5) -- (v20) -- cycle;
\fill[blue!3] (v2) -- (v7) -- (v13) -- cycle;
\fill[blue!3] (v2) -- (v7) -- (v20) -- cycle;
\fill[blue!3] (v2) -- (v10) -- (v14) -- cycle;
\fill[blue!3] (v2) -- (v13) -- (v14) -- cycle;
\fill[blue!3] (v3) -- (v6) -- (v16) -- cycle;
\fill[blue!3] (v3) -- (v6) -- (v20) -- cycle;
\fill[blue!3] (v3) -- (v7) -- (v19) -- cycle;
\fill[blue!3] (v3) -- (v7) -- (v20) -- cycle;
\fill[blue!3] (v3) -- (v8) -- (v16) -- cycle;
\fill[blue!3] (v3) -- (v8) -- (v19) -- cycle;
\fill[blue!3] (v4) -- (v9) -- (v15) -- cycle;
\fill[blue!3] (v5) -- (v10) -- (v18) -- cycle;
\fill[blue!3] (v6) -- (v12) -- (v16) -- cycle;
\fill[blue!3] (v7) -- (v13) -- (v19) -- cycle;
\draw[draw=black!52,line width=.38pt] (v0) -- (v4);
\draw[draw=black!52,line width=.38pt] (v0) -- (v5);
\draw[draw=black!52,line width=.38pt] (v0) -- (v11);
\draw[draw=black!52,line width=.38pt] (v0) -- (v15);
\draw[draw=black!52,line width=.38pt] (v0) -- (v18);
\draw[draw=black!52,line width=.38pt] (v0) -- (v20);
\draw[draw=black!52,line width=.38pt] (v1) -- (v4);
\draw[draw=black!52,line width=.38pt] (v1) -- (v6);
\draw[draw=black!52,line width=.38pt] (v1) -- (v9);
\draw[draw=black!52,line width=.38pt] (v1) -- (v12);
\draw[draw=black!52,line width=.38pt] (v1) -- (v17);
\draw[draw=black!52,line width=.38pt] (v1) -- (v20);
\draw[draw=black!52,line width=.38pt] (v2) -- (v5);
\draw[draw=black!52,line width=.38pt] (v2) -- (v7);
\draw[draw=black!52,line width=.38pt] (v2) -- (v10);
\draw[draw=black!52,line width=.38pt] (v2) -- (v13);
\draw[draw=black!52,line width=.38pt] (v2) -- (v14);
\draw[draw=black!52,line width=.38pt] (v2) -- (v20);
\draw[draw=black!52,line width=.38pt] (v3) -- (v6);
\draw[draw=black!52,line width=.38pt] (v3) -- (v7);
\draw[draw=black!52,line width=.38pt] (v3) -- (v8);
\draw[draw=black!52,line width=.38pt] (v3) -- (v16);
\draw[draw=black!52,line width=.38pt] (v3) -- (v19);
\draw[draw=black!52,line width=.38pt] (v3) -- (v20);
\draw[draw=black!52,line width=.38pt] (v4) -- (v9);
\draw[draw=black!52,line width=.38pt] (v4) -- (v15);
\draw[draw=black!52,line width=.38pt] (v4) -- (v20);
\draw[draw=black!52,line width=.38pt] (v5) -- (v10);
\draw[draw=black!52,line width=.38pt] (v5) -- (v18);
\draw[draw=black!52,line width=.38pt] (v5) -- (v20);
\draw[draw=black!52,line width=.38pt] (v6) -- (v12);
\draw[draw=black!52,line width=.38pt] (v6) -- (v16);
\draw[draw=black!52,line width=.38pt] (v6) -- (v20);
\draw[draw=black!52,line width=.38pt] (v7) -- (v13);
\draw[draw=black!52,line width=.38pt] (v7) -- (v19);
\draw[draw=black!52,line width=.38pt] (v7) -- (v20);
\draw[draw=witnessred,line width=.8pt] (v8) -- (v16);
\draw[draw=witnessred,line width=.8pt] (v8) -- (v19);
\draw[draw=witnessred,line width=.8pt] (v9) -- (v15);
\draw[draw=witnessred,line width=.8pt] (v9) -- (v17);
\draw[draw=witnessred,line width=.8pt] (v10) -- (v14);
\draw[draw=witnessred,line width=.8pt] (v10) -- (v18);
\draw[draw=witnessred,line width=.8pt] (v11) -- (v15);
\draw[draw=witnessred,line width=.8pt] (v11) -- (v18);
\draw[draw=witnessred,line width=.8pt] (v12) -- (v16);
\draw[draw=witnessred,line width=.8pt] (v12) -- (v17);
\draw[draw=witnessred,line width=.8pt] (v13) -- (v14);
\draw[draw=witnessred,line width=.8pt] (v13) -- (v19);
\node[jfef vertex] at (v0) {$0$};
\node[jfef vertex] at (v1) {$1$};
\node[jfef vertex] at (v2) {$2$};
\node[jfef vertex] at (v3) {$3$};
\node[jfef vertex] at (v4) {$4$};
\node[jfef vertex] at (v5) {$5$};
\node[jfef vertex] at (v6) {$6$};
\node[jfef vertex] at (v7) {$7$};
\node[jfef vertex,draw=witnessred,line width=.7pt] at (v8) {$8$};
\node[jfef vertex,draw=witnessred,line width=.7pt] at (v9) {$8$};
\node[jfef vertex,draw=witnessred,line width=.7pt] at (v10) {$8$};
\node[jfef vertex,draw=witnessred,line width=.7pt] at (v11) {$9$};
\node[jfef vertex,draw=witnessred,line width=.7pt] at (v12) {$9$};
\node[jfef vertex,draw=witnessred,line width=.7pt] at (v13) {$9$};
\node[jfef vertex,draw=witnessred,line width=.7pt] at (v14) {$10$};
\node[jfef vertex,draw=witnessred,line width=.7pt] at (v15) {$10$};
\node[jfef vertex,draw=witnessred,line width=.7pt] at (v16) {$10$};
\node[jfef vertex,draw=witnessred,line width=.7pt] at (v17) {$11$};
\node[jfef vertex,draw=witnessred,line width=.7pt] at (v18) {$11$};
\node[jfef vertex,draw=witnessred,line width=.7pt] at (v19) {$11$};
\node[jfef vertex] at (v20) {$12$};
\end{tikzpicture}
\hspace{6mm}\begin{tikzpicture}[x=1cm,y=1cm]
\definecolor{witnessred}{RGB}{181,55,36}
\definecolor{witnessblue}{RGB}{32,89,140}
\coordinate (v0) at (1.146661411,2.139452559);
\coordinate (v1) at (-2.265488516,1.039851549);
\coordinate (v2) at (2.265489343,-1.039852163);
\coordinate (v3) at (-1.146663831,-2.139454358);
\coordinate (v4) at (-1.146661411,2.139452559);
\coordinate (v5) at (2.265488516,1.039851549);
\coordinate (v6) at (-2.265489343,-1.039852163);
\coordinate (v7) at (1.146663831,-2.139454358);
\coordinate (v8) at (3.440000000,-3.200000000);
\coordinate (v9) at (-3.440000000,3.200000000);
\coordinate (v10) at (3.440000000,3.200000000);
\coordinate (v11) at (-3.440000000,-3.200000000);
\coordinate (v12) at (0.000000000,0.000000000);
\draw[draw=black!52,line width=.38pt] (v0) -- (v4);
\draw[draw=black!52,line width=.38pt] (v0) -- (v5);
\draw[draw=black!52,line width=.38pt] (v0) -- (v9);
\draw[draw=black!52,line width=.38pt] (v0) -- (v10);
\draw[draw=black!52,line width=.38pt] (v0) -- (v11);
\draw[draw=black!52,line width=.38pt] (v0) -- (v12);
\draw[draw=black!52,line width=.38pt] (v1) -- (v4);
\draw[draw=black!52,line width=.38pt] (v1) -- (v6);
\draw[draw=black!52,line width=.38pt] (v1) -- (v8);
\draw[draw=black!52,line width=.38pt] (v1) -- (v9);
\draw[draw=black!52,line width=.38pt] (v1) -- (v11);
\draw[draw=black!52,line width=.38pt] (v1) -- (v12);
\draw[draw=black!52,line width=.38pt] (v2) -- (v5);
\draw[draw=black!52,line width=.38pt] (v2) -- (v7);
\draw[draw=black!52,line width=.38pt] (v2) -- (v8);
\draw[draw=black!52,line width=.38pt] (v2) -- (v9);
\draw[draw=black!52,line width=.38pt] (v2) -- (v10);
\draw[draw=black!52,line width=.38pt] (v2) -- (v12);
\draw[draw=black!52,line width=.38pt] (v3) -- (v6);
\draw[draw=black!52,line width=.38pt] (v3) -- (v7);
\draw[draw=black!52,line width=.38pt] (v3) -- (v8);
\draw[draw=black!52,line width=.38pt] (v3) -- (v10);
\draw[draw=black!52,line width=.38pt] (v3) -- (v11);
\draw[draw=black!52,line width=.38pt] (v3) -- (v12);
\draw[draw=black!52,line width=.38pt] (v4) -- (v8);
\draw[draw=black!52,line width=.38pt] (v4) -- (v10);
\draw[draw=black!52,line width=.38pt] (v4) -- (v12);
\draw[draw=black!52,line width=.38pt] (v5) -- (v8);
\draw[draw=black!52,line width=.38pt] (v5) -- (v11);
\draw[draw=black!52,line width=.38pt] (v5) -- (v12);
\draw[draw=black!52,line width=.38pt] (v6) -- (v9);
\draw[draw=black!52,line width=.38pt] (v6) -- (v10);
\draw[draw=black!52,line width=.38pt] (v6) -- (v12);
\draw[draw=black!52,line width=.38pt] (v7) -- (v9);
\draw[draw=black!52,line width=.38pt] (v7) -- (v11);
\draw[draw=black!52,line width=.38pt] (v7) -- (v12);
\draw[draw=witnessred,line width=.8pt] (v8) -- (v10);
\draw[draw=witnessred,line width=.8pt] (v8) -- (v11);
\draw[draw=witnessred,line width=.8pt] (v9) -- (v10);
\draw[draw=witnessred,line width=.8pt] (v9) -- (v11);
\node[jfef vertex] at (v0) {$0$};
\node[jfef vertex] at (v1) {$1$};
\node[jfef vertex] at (v2) {$2$};
\node[jfef vertex] at (v3) {$3$};
\node[jfef vertex] at (v4) {$4$};
\node[jfef vertex] at (v5) {$5$};
\node[jfef vertex] at (v6) {$6$};
\node[jfef vertex] at (v7) {$7$};
\node[jfef vertex,draw=witnessred,line width=.7pt] at (v8) {$8$};
\node[jfef vertex,draw=witnessred,line width=.7pt] at (v9) {$9$};
\node[jfef vertex,draw=witnessred,line width=.7pt] at (v10) {$10$};
\node[jfef vertex,draw=witnessred,line width=.7pt] at (v11) {$11$};
\node[jfef vertex] at (v12) {$12$};
\end{tikzpicture}

\caption{On the left, a triangulated disc whose equally labelled boundary
vertices and corresponding boundary edges are identified to obtain
$\cl(\Gt)$. The boundary traverses
$10$--$9$--$11$--$8$ three times. The $28$ small triangular regions are
its faces. On the right, the graph $\Gt$; the four edges in three
triangles are highlighted. The highlighted boundary on the left maps
to the highlighted four-cycle on the right.}\label{fig:G13}
\end{figure}

The construction also attains the least possible order for odd torsion
in the first homology of a flag complex. This conclusion uses the
earlier classification of torsion graphs, without an acyclicity
assumption.

\begin{proposition}\label{rem:adamaszek}
The minimum number of vertices of a flag complex whose first integral
homology has odd torsion is $13$.
\end{proposition}

\begin{proof}
A torsion graph is called irreducible if every vertex deletion removes
all torsion from its first homology. Adamaszek's classification
\cite[Theorem~1.1 and Section~4]{Adamaszek2014} gives no torsion through
ten vertices, four torsion graphs on eleven vertices, each with first
homology $\ZZ/2$, and $363$ irreducible torsion graphs on twelve vertices,
each with torsion subgroup $\ZZ/2$. The absence of torsion through ten vertices and the four torsion
graphs on eleven vertices also follow from Katzman
\cite[Theorem~4.1 and Section~5]{Katzman2006}.

If a twelve-vertex torsion graph $G$ is not irreducible, it has a vertex
$v$ such that $G-v$ is one of the four eleven-vertex graphs. If $v$ is
isolated, then $H_1(\cl(G);\ZZ)=H_1(\cl(G-v);\ZZ)\cong\ZZ/2$.
Otherwise, apply the Mayer--Vietoris sequence to
\[
 \cl(G)=\cl(G-v)\cup\bigl(v*\cl(\Nb_G(v))\bigr).
\]
It gives an exact sequence
$0\to Q\to H_1(\cl(G);\ZZ)\to F\to0$, where $Q$ is a quotient of
$\ZZ/2$ and $F$ is a subgroup of
$\widetilde H_0(\cl(\Nb_G(v));\ZZ)$. The group $F$ is free abelian,
so the sequence splits. Thus $H_1(\cl(G);\ZZ)$ has no odd torsion.
Theorem~\ref{thm:A} attains the lower bound.
\end{proof}

\subsection{A non-strong I-contractible graph on fourteen vertices}
\label{sec:fourteen}

For the second assertion of Conjecture~14, nonzero homology cannot
provide the obstruction: an I-contractible graph is integrally
acyclic. Instead, we show that every possible first strong deletion
is obstructed, and then establish I-contractibility by an edge addition.

\begin{theorem}\label{thm:B}
The graph $\Ga$ is $K_4$-free, has $14$ vertices and $44$ edges, and
belongs to $\I\setminus\IS$. Its one-edge extension $\Ga+\{0,1\}$ belongs to
$\IS$, and deletion of that edge is $\IS$-admissible.
Consequently, $\Ga$ refutes the second assertion of Conjecture~14,
and $\Ga+\{0,1\}$ refutes Conjecture~15.
\end{theorem}

\begin{proof}
The common neighbourhood of $0$ and $1$ in $\Ga$ is the induced path
\[
 8\text{--}12\text{--}5\text{--}10\text{--}13.
\]
Since nonempty trees belong to $\IS$ (Section~\ref{sec:contractible}),
adding $\{0,1\}$ is admissible.
Starting from $\Ga+\{0,1\}$, perform the following deletions in the
displayed order:
\[
\begin{gathered}
 e(0,5),\ e(0,10),\ e(0,12),\ e(8,12),\ e(1,12),\\
 v(5),\ v(10),\ v(3),\ v(6),\ v(2),\ v(7),\ v(1),\ v(0),\\
 v(4),\ v(8),\ v(11),\ v(12),\ v(9).
\end{gathered}
\]
Here $e(u,w)$ deletes an edge and $v(u)$ deletes a vertex; labels do
not change after a deletion. Every tested neighbourhood is a nonempty
tree, as the explicit lists in Table~\ref{tab:Ga-reduction} show.
The final graph is the singleton $13$. Hence $\Ga+\{0,1\}\in\IS$.
Deleting the added edge is an I-admissible transformation, with the
same common neighbourhood as before its addition, so $\Ga\in\I$.

To prove $\Ga\notin\IS$, consider its links. Each vertex neighbourhood
is a connected triangle-free graph; its first Betti number is positive
in every case, as recorded in Table~\ref{tab:witness-links}.
Every edge has two or three common neighbours, and those neighbours
form an independent set because $\Ga$ is $K_4$-free.
Thus neither a vertex neighbourhood nor an edge neighbourhood is
acyclic. None belongs to $\IS$, so no first strong deletion is possible.

The two graphs have $14$ vertices, whereas both examples conjectured
to be smallest in the earlier paper have $15$ vertices. This proves
the stated refutations. The global lower bounds are established from
the complete census below.
\end{proof}

\begin{figure}[tbp]
\centering
\begin{tikzpicture}[x=1cm,y=1cm]
\definecolor{witnessred}{RGB}{181,55,36}
\definecolor{witnessblue}{RGB}{32,89,140}
\coordinate (v0) at (4.899000000,-0.124598525);
\coordinate (v1) at (2.722969139,3.944383940);
\coordinate (v2) at (-2.435697872,0.406103749);
\coordinate (v3) at (-4.899000000,4.045228676);
\coordinate (v4) at (1.271029059,-0.654558239);
\coordinate (v5) at (4.387943889,-3.466555902);
\coordinate (v6) at (-4.899000000,-4.554000000);
\coordinate (v7) at (-2.939555778,-3.846959139);
\coordinate (v8) at (-0.868699866,3.253825943);
\coordinate (v9) at (2.353335221,-2.606664647);
\coordinate (v10) at (3.219437385,-4.554000000);
\coordinate (v11) at (3.996206901,3.171731296);
\coordinate (v12) at (1.198003432,4.416000000);
\coordinate (v13) at (-1.343328314,-3.846959139);
\draw[draw=witnessblue,line width=1pt,dashed] (v0) -- (v1);
\draw[draw=black!52,line width=.38pt] (v0) -- (v4);
\draw[draw=black!52,line width=.38pt] (v0) -- (v5);
\draw[draw=black!52,line width=.38pt] (v0) -- (v8);
\draw[draw=black!52,line width=.38pt] (v0) -- (v10);
\draw[draw=black!52,line width=.38pt] (v0) -- (v12);
\draw[draw=black!52,line width=.38pt] (v0) -- (v13);
\draw[draw=black!52,line width=.38pt] (v1) -- (v5);
\draw[draw=black!52,line width=.38pt] (v1) -- (v7);
\draw[draw=black!52,line width=.38pt] (v1) -- (v8);
\draw[draw=black!52,line width=.38pt] (v1) -- (v10);
\draw[draw=black!52,line width=.38pt] (v1) -- (v12);
\draw[draw=black!52,line width=.38pt] (v1) -- (v13);
\draw[draw=black!52,line width=.38pt] (v2) -- (v6);
\draw[draw=black!52,line width=.38pt] (v2) -- (v7);
\draw[draw=black!52,line width=.38pt] (v2) -- (v9);
\draw[draw=black!52,line width=.38pt] (v2) -- (v12);
\draw[draw=black!52,line width=.38pt] (v2) -- (v13);
\draw[draw=black!52,line width=.38pt] (v3) -- (v6);
\draw[draw=black!52,line width=.38pt] (v3) -- (v9);
\draw[draw=black!52,line width=.38pt] (v3) -- (v10);
\draw[draw=black!52,line width=.38pt] (v3) -- (v12);
\draw[draw=black!52,line width=.38pt] (v3) -- (v13);
\draw[draw=black!52,line width=.38pt] (v4) -- (v8);
\draw[draw=black!52,line width=.38pt] (v4) -- (v9);
\draw[draw=black!52,line width=.38pt] (v4) -- (v11);
\draw[draw=black!52,line width=.38pt] (v4) -- (v13);
\draw[draw=black!52,line width=.38pt] (v5) -- (v9);
\draw[draw=witnessred,line width=.8pt] (v5) -- (v10);
\draw[draw=black!52,line width=.38pt] (v5) -- (v11);
\draw[draw=witnessred,line width=.8pt] (v5) -- (v12);
\draw[draw=black!52,line width=.38pt] (v6) -- (v7);
\draw[draw=black!52,line width=.38pt] (v6) -- (v10);
\draw[draw=black!52,line width=.38pt] (v6) -- (v11);
\draw[draw=black!52,line width=.38pt] (v6) -- (v12);
\draw[draw=black!52,line width=.38pt] (v7) -- (v8);
\draw[draw=black!52,line width=.38pt] (v7) -- (v11);
\draw[draw=black!52,line width=.38pt] (v7) -- (v13);
\draw[draw=black!52,line width=.38pt] (v8) -- (v11);
\draw[draw=witnessred,line width=.8pt] (v8) -- (v12);
\draw[draw=black!52,line width=.38pt] (v9) -- (v11);
\draw[draw=black!52,line width=.38pt] (v9) -- (v12);
\draw[draw=black!52,line width=.38pt] (v9) -- (v13);
\draw[draw=black!52,line width=.38pt] (v10) -- (v11);
\draw[draw=witnessred,line width=.8pt] (v10) -- (v13);
\node[jfef vertex,draw=witnessblue,line width=.8pt] at (v0) {$0$};
\node[jfef vertex,draw=witnessblue,line width=.8pt] at (v1) {$1$};
\node[jfef vertex] at (v2) {$2$};
\node[jfef vertex] at (v3) {$3$};
\node[jfef vertex] at (v4) {$4$};
\node[jfef vertex,fill=orange!12] at (v5) {$5$};
\node[jfef vertex] at (v6) {$6$};
\node[jfef vertex] at (v7) {$7$};
\node[jfef vertex,fill=orange!12] at (v8) {$8$};
\node[jfef vertex] at (v9) {$9$};
\node[jfef vertex,fill=orange!12] at (v10) {$10$};
\node[jfef vertex] at (v11) {$11$};
\node[jfef vertex,fill=orange!12] at (v12) {$12$};
\node[jfef vertex,fill=orange!12] at (v13) {$13$};
\end{tikzpicture}

\caption{The graph $\Ga$ consists of the solid edges. Adding the dashed
edge $\{0,1\}$ gives a strong I-contractible graph. The common
neighbourhood of its endpoints is the highlighted induced path
$8$--$12$--$5$--$10$--$13$. Deleting the dashed edge is admissible but
leaves a graph outside $\IS$.}\label{fig:G14}
\end{figure}

The order-$14$ classification also uses $\Gb$ and $\Gc$, two further
members of $\I\setminus\IS$ that become strong I-contractible after
one admissible edge addition. Table~\ref{tab:witness-adjacency} defines
them and Table~\ref{tab:variant-reductions} supplies their reductions.

\begin{proposition}\label{prop:other-seeds}
The graphs $\Gb$ and $\Gc$ are $K_4$-free members of
$\I\setminus\IS$, each with $14$ vertices and $45$ edges.
The extensions $\Gb+\{0,3\}$ and $\Gc+\{2,3\}$ belong to $\IS$,
and the added edges are admissible.
\end{proposition}

\begin{proof}
The common neighbourhoods for the two additions are the induced paths
$7\text{--}13\text{--}9\text{--}12$ and
$6\text{--}13\text{--}7\text{--}10\text{--}11$, respectively.
In each column of Table~\ref{tab:variant-reductions}, every tested
neighbourhood is a nonempty tree and the last vertex is $13$.
This proves membership of both extensions in $\IS$ and of the original
graphs in $\I$. Table~\ref{tab:witness-links} excludes every possible
first strong deletion from $\Gb$ and $\Gc$, just as for $\Ga$.
The graph $\Gb$ has an edge with four common neighbours, whereas
$\Gc$ has none. Together with their edge counts, this shows that
$\Ga,\Gb,\Gc$ are pairwise non-isomorphic.
\end{proof}

The same local structure also obstructs simplicial collapses.

\begin{proposition}\label{prop:noncollapsible}
The clique complexes of $\Ga,\Gb,\Gc$ are contractible and
non-collapsible.
\end{proposition}

\begin{proof}
Their I-contractibility implies contractibility of their clique
complexes. Each complex is two-dimensional. By
Table~\ref{tab:witness-links}, every edge lies in at least two triangles,
and every vertex lies in a triangle. There is therefore no free face
with which to start an elementary collapse.
\end{proof}

\section{Algorithms and the enumeration of graph classes}\label{sec:algorithms}

The table requires decisions about a graph and about the neighbourhoods
that occur during its reductions. We first give exact procedures for
these decisions. We then describe how to enumerate isomorphism classes
and how to reduce the classification at order $14$ to a smaller family.

\subsection{Cliques and homology}

Fix an ordering of the vertices of a nonempty graph $G$. Starting with
each vertex, extend an increasing clique by adjoining a larger vertex
adjacent to all its members, recording each clique as it is formed.
Continue until no extension is possible.
Every clique has a unique increasing order and a unique sequence of
initial segments, so this procedure lists it exactly once. In particular,
it determines all dimensions of $\cl(G)$, not only its triangles.

For an increasing clique $[v_0,\ldots,v_i]$, use the boundary
\[
 \partial_i[v_0,\ldots,v_i]
   =\sum_{j=0}^i(-1)^j[v_0,\ldots,\widehat v_j,\ldots,v_i].
\]
Adjoin the augmentation $\partial_0:C_0\to\ZZ$, which sends every
vertex to $1$. Write $f_i=\operatorname{rank} C_i$ and let $r_i(R)$
denote the rank of $\partial_i$ over the field $R$. Since consecutive
boundaries compose to zero, elementary linear algebra gives
\begin{equation}\label{eq:betti}
 \dim_R\widetilde H_i(G;R)=f_i-r_i(R)-r_{i+1}(R)
 \qquad(i\ge0),
\end{equation}
where $r_0(R)=1$ and the rank above the highest dimension is zero.
Thus elimination over $\mathbb F_2$ decides membership in
$\mathcal A_{\mathbb F_2}$.
The equality $\chi(G)=1$ is a necessary preliminary condition, but is
not a substitute for the homology calculation.

Integral homology also requires torsion information. If $d_{i+1,j}$
are the nonzero Smith factors of $\partial_{i+1}$, then
\begin{equation}\label{eq:integral-homology}
 \widetilde H_i(G;\ZZ)\cong
 \ZZ^{\,f_i-r_i(\mathbb Q)-r_{i+1}(\mathbb Q)}
 \oplus\bigoplus_{d_{i+1,j}>1}\ZZ/d_{i+1,j}.
\end{equation}
Indeed, the augmented chain complex gives a short exact sequence
\[
 0\longrightarrow\widetilde H_i(G;\ZZ)
 \longrightarrow C_i/\operatorname{im}\partial_{i+1}
 \longrightarrow\operatorname{im}\partial_i\longrightarrow0.
\]
Its last term is free abelian, so the sequence splits. The Smith
factors describe the middle term, and the two adjacent rational ranks
give the displayed free rank. Formula~\eqref{eq:integral-homology},
in every dimension present, therefore decides integral acyclicity.

\subsection{Strong and vertex reductions}

Let $S(G)$ be a truth value. Set $S(\emptyset)=\mathrm{false}$ and
$S(K_1)=\mathrm{true}$. For a graph with at least two vertices evaluate
\begin{align}
 S(G)={}&\bigvee_{v\in V(G)}
 \bigl(S(\Nb_G(v))\land S(G-v)\bigr)\notag\\
 &{}\lor\bigvee_{uv\in E(G)}
 \bigl(S(\Nb_G(u,v))\land S(G-uv)\bigr).
 \label{eq:strong-rec}
\end{align}
Evaluation may stop when one conjunction is true; a negative answer
requires every conjunction to be false. The argument of $S$ is always
the current graph, including the edges that remain after earlier
deletions.

\begin{proposition}\label{prop:IS-correct}
The evaluation of~\eqref{eq:strong-rec} terminates and returns true
exactly when $G\in\IS$. A true answer determines a reduction to $K_1$
and reductions of the neighbourhoods that justify its steps.
\end{proposition}
\begin{proof}
Induct on $|V(G)|+|E(G)|$. Every neighbourhood and every remaining
graph in~\eqref{eq:strong-rec} has smaller measure, and there are
finitely many alternatives. This proves termination, starting from
the two stated bases. If a conjunction is true, induction supplies
the two reductions that make the corresponding first deletion valid;
prepending that deletion proves $G\in\IS$. Conversely, a strong
reduction of a non-singleton graph has a first vertex or edge deletion.
Its neighbourhood and remainder are in $\IS$, so induction makes
the corresponding conjunction true. Recording a successful choice
at each occurrence gives the asserted reductions.
\end{proof}

For vertex reductions, all intermediate graphs are induced subgraphs
of a fixed graph $G$. For $U\subseteq V(G)$ and $v\in U$, put
$N_U(v)=U\cap V(\Nb_G(v))$. Define truth values $W(U)$ and $D_k(U)$
to be false on the empty set and true on singletons. For nonempty $U$,
$D_{-1}(U)$ means that $G[U]$ is a cone. On sets of size at least two,
use
\begin{align}
 W(U)&=\bigvee_{v\in U}
     \bigl(W(N_U(v))\land W(U\setminus\{v\})\bigr),
     \label{eq:ivs-rec}\\
 D_k(U)&=\bigvee_{v\in U}
     \bigl(D_{k-1}(N_U(v))\land D_k(U\setminus\{v\})\bigr),
     \qquad k\ge0.\label{eq:kdis-rec}
\end{align}

\begin{proposition}\label{prop:vertex-correct}
The recurrences~\eqref{eq:ivs-rec} and~\eqref{eq:kdis-rec} decide
whether $G[U]$ belongs to $\ISV$ and to $k$-Dis, respectively. For a
fixed largest level, evaluation in increasing order of $|U|$
terminates.
\end{proposition}
\begin{proof}
Each set on the right is a proper subset of $U$. Induct on its
cardinality, evaluating the cone test directly. A successful
conjunction supplies a first admissible vertex and a reduction of
the remainder. In the opposite direction, the first step of any
reduction supplies one of the conjunctions. The neighbourhood
condition is at the same recursive class for $\ISV$ and at level
$k-1$ for $k$-Dis, exactly as in their definitions. This proves both
equivalences and termination.
\end{proof}

At level zero there is also a simpler procedure: delete any dominated
vertex until no such vertex remains. Here $v$ is dominated by $w\ne v$
when $N_G[v]\subseteq N_G[w]$, where $N_G[v]=\{v\}\cup V(\Nb_G(v))$.
This is equivalent to $\Nb_G(v)$ being a cone with apex $w$.
The procedure succeeds precisely on the $0$-dismantlable graphs
\cite[Section~4.1]{FieuxJouve2020}. To see why the order does not
matter, consider two available deletions. Unless the vertices dominate
each other, each remains dominated after the other is deleted:
if a dominator is removed, use its dominator and the transitivity of
closed-neighbourhood inclusion. Mutual domination gives identical
closed neighbourhoods, so the two remainders are isomorphic.
Induction on the number of vertices now shows that an arbitrary
available deletion preserves dismantlability. Reaching $K_1$ is itself
a dismantling, proving the converse. This procedure gives the
$0$-Dis column of the table.
No analogous order-independence holds for~\eqref{eq:strong-rec}:
the extensions in Section~\ref{sec:examples} admit a deletion that
leaves a non-strong graph.

\subsection{Positive tests for I-contractibility}

Strong reductions also supply positive tests for the larger class
$\I$. Suppose
\begin{equation}\label{eq:I-path}
 G=G_0,\qquad G_j=G_{j-1}+u_jv_j\quad(1\le j\le t),
 \qquad G_t\in\IS,
\end{equation}
and every added edge has common neighbourhood in $\IS$ at its step.
Then $G\in\I$. Indeed, start from $G_t\in\IS\subseteq\I$ and delete
the added edges in reverse order. Adding or deleting the edge between
two vertices leaves their common neighbourhood unchanged, so each
reverse step is an I-transformation.

For a fixed bound $t$, enumerate all admissible edge additions from
$G$, one step at a time, and test each reached graph by
Proposition~\ref{prop:IS-correct}. This finds a path of the form
\eqref{eq:I-path} of length at most $t$ whenever one exists: induction
on the path length shows that every possible successor is considered.
The search is finite, since the vertex set is fixed and edges are
only added. Its failure excludes only the specified paths, not
membership in $\I$. The negative decisions needed below instead use
nonzero integral homology. Together these two tests will settle all
graphs left outside $\IS$ in the universes of the table.

\subsection{Representatives up to isomorphism}\label{sec:representatives}

The generation uses canonical vertex augmentation
\cite[Section~4]{McKay1998}. For each graph $H$ of order at least
two, choose a vertex orbit $O(H)$ under $\operatorname{Aut}(H)$,
invariantly under isomorphism. This choice can be made from a
canonical labelled representative. Begin with $K_1$. For every
retained graph $G$, adjoin a vertex $x$ adjacent to one subset from
each orbit of subsets of $V(G)$ under $\operatorname{Aut}(G)$.
Retain the resulting graph $H$ when $x\in O(H)$. Connectedness is
required at the output order. For the $K_4$-free universe, reject
augmentations containing a $K_4$.

\begin{proposition}\label{prop:enumeration-correct}
This construction, followed by the connectedness test, generates
exactly one representative of every connected graph of the requested
order. The additional $K_4$ rejection gives exactly the connected
$K_4$-free graphs.
\end{proposition}
\begin{proof}
Induct on the order, before testing connectedness. For a
target $H$, choose $x\in O(H)$. Its parent is represented by induction
and the neighbourhood orbit is considered. The resulting copy of
$H$ is accepted because $O$ is invariant under isomorphism, proving
existence.
For two isomorphic accepted augmentations, an isomorphism followed
by an automorphism identifies the new vertices, and hence the parents
and neighbourhood orbits. These were considered once, proving
uniqueness. Restricting the outputs preserves both conclusions.
Every induced subgraph of a $K_4$-free graph is $K_4$-free, so the
additional rejection loses no predecessor of an eligible target.
\end{proof}

For any of the exact predicates above, sum its value over the generated
representatives of a fixed order. Proposition~\ref{prop:enumeration-correct}
and the decision propositions show that this sum is the cardinality
of the corresponding class, with neither missing isomorphism classes
nor repeated contributions. The column for $\I$ additionally requires
the positive and negative adjudications described in
Section~\ref{sec:table}.
The inclusions in~\eqref{eq:dismantling-inclusions} and
\eqref{eq:homology-inclusions} allow a positive decision for one
class to establish membership in every larger class of these chains;
a negative decision establishes non-membership in every smaller class.

\subsection{Reduction of the order-fourteen classification}\label{sec:reduction}

The complete census gives the class counts. To identify the acyclic
non-strong graphs behind these counts, we reduce them until no
admissible move remains.

\begin{definition}\label{def:stuck}
A graph of order at least two is \defi{stuck} if it is
$\mathbb F_2$-acyclic and has no $\IS$-admissible vertex or edge.
A graph of order at least two is \defi{almost stuck at $u$} if it is
$\mathbb F_2$-acyclic,
has no $\IS$-admissible edge, and no vertex other than $u$ is
$\IS$-admissible. The distinguished vertex $u$ need not be admissible.
\end{definition}

Every stuck graph lies outside $\IS$, since a reduction cannot start.
Two observations explain the role of minimum order. An admissible
deletion from a graph of $\I$ remains in $\I$. Moreover, if its
remainder is in $\IS$, then the original graph is in $\IS$, by
prepending that deletion to a strong reduction.

\begin{lemma}\label{lem:stuck}
Let $\mathcal H$ be closed under vertex and edge deletion, and let $n$
be the smallest order attained in $\mathcal H\cap(\I\setminus\IS)$.
Every graph in this class of order $n$ reduces by admissible edge
deletions to a stuck graph in the same class.
\end{lemma}
\begin{proof}
An admissible vertex deletion would give a smaller graph in the same
class, by the two observations, contradicting the choice of $n$.
Delete admissible edges as long as possible. The order is unchanged,
membership in $\mathcal H\cap(\I\setminus\IS)$ is preserved, and
the number of edges decreases. At termination neither kind of
deletion is possible. The terminal graph is acyclic because it is
in $\I$, so it is stuck.
\end{proof}

\begin{lemma}\label{lem:2conn}
If a stuck graph has a cut vertex $u$, it is the union of two induced
graphs $A,B$, each of order at least three, that meet only at $u$ and
are almost stuck there. Consequently, if no almost stuck graph in a
hereditary class has order at most $m$, every stuck graph in that
class of order at most $2m$ is $2$-connected.
\end{lemma}
\begin{proof}
Partition the components of $G-u$ into two nonempty collections and
include $u$ in both pieces. Their clique complexes meet at one
vertex, so reduced Mayer--Vietoris gives
\[
 \widetilde H_i(G;\mathbb F_2)\cong
 \widetilde H_i(A;\mathbb F_2)\oplus
 \widetilde H_i(B;\mathbb F_2).
\]
Both pieces are acyclic. The neighbourhood of a vertex other than
$u$, and the common neighbourhood of every edge in a piece, agree
with those in $G$. Thus no such element is admissible. A piece of
order two would have a leaf of $G$, which is admissible. Both orders
are therefore at least three. If both exceed $m$, their union has
order at least $2(m+1)-1=2m+1$, proving the last assertion.
\end{proof}

\begin{lemma}\label{lem:degrees}
Every edge of a $2$-connected stuck graph $G$ lies in at least two
triangles, and
\[
 4\le\delta(G)\le\Delta(G)\le |V(G)|-2.
\]
\end{lemma}
\begin{proof}
An edge in no triangle belongs to a cycle, by $2$-connectedness.
Its coefficient in every $2$-boundary is zero, whereas its coefficient
in that cycle is one. The cycle therefore represents nonzero homology
over $\mathbb F_2$, contradicting acyclicity.
An edge in exactly one triangle has common neighbourhood $K_1$ and
is admissible. Hence every edge lies in at least two triangles.
For $H=\Nb_G(v)$, every vertex $w$ of $H$ consequently has at least
two neighbours in $H$. If $|V(H)|\le3$, this forces $H=K_3$, making
$v$ admissible. Thus every degree is at least four. A universal
vertex would make $G$ a cone and hence a member of $\IS$, so no
degree is $|V(G)|-1$.
\end{proof}

On $K_4$-free graphs these reductions have particularly simple tests.
\begin{proposition}\label{rem:K4-local}
In a $K_4$-free graph, a vertex is $\IS$-admissible exactly when its
open neighbourhood is a nonempty tree, and an edge is $\IS$-admissible
exactly when it belongs to one triangle.
\end{proposition}
\begin{proof}
A vertex neighbourhood is triangle-free. A triangle-free graph in
$\IS$ has a contractible clique complex of dimension at most one,
so it is a tree. Conversely, a nonempty tree reduces by leaf deletions
to $K_1$. The common neighbourhood of an edge is independent and
belongs to $\IS$ exactly when it has one vertex. Its vertices
correspond to the triangles containing the edge.
\end{proof}

This criterion is independent of any small-order census. In particular,
to test stuckness on a $K_4$-free graph, check its reduced homology and
exclude every tree vertex neighbourhood and every edge in one triangle.
For almost stuckness allow at most one such vertex, including the
possibility of none.

Finally, reverse the edge deletions of Lemma~\ref{lem:stuck} to recover
the graphs arising from a stuck graph. Starting with a finite set of
stuck graphs, called \defi{seeds}, adjoin every missing edge with common neighbourhood in
$\IS$; continue from non-strong successors and discard strong ones.
Identify isomorphic graphs throughout. The retained set is the
\defi{pruned edge-addition closure}.

\begin{proposition}\label{prop:closure-correct}
This procedure terminates and retains exactly the non-strong graphs
reachable from the seeds by admissible edge additions. To recover
only $K_4$-free targets, one may also discard successors containing
a $K_4$.
\end{proposition}
\begin{proof}
There are finitely many graphs on the fixed vertex set and every
step adds an edge. All retained graphs are reachable, by construction.
If $H\in\IS$ and $H+uv$ is an admissible addition, deleting $uv$
first and then reducing $H$ proves $H+uv\in\IS$. A path to a
non-strong target therefore never passes through a strong graph.
Every predecessor on that path is expanded, so induction on path
length proves completeness. Isomorphisms transport nonedges and
their neighbourhoods, and hence all possible continuations.
Finally, adding edges cannot remove a $K_4$, so no path to a
$K_4$-free target is lost by the additional restriction.
\end{proof}

\section{Cardinalities and minimum orders}\label{sec:searches}\label{sec:table}

We count connected graphs up to isomorphism. For orders $3$ through
$13$ the universe is unrestricted; at order $14$ it consists of the
connected $K_4$-free graphs. The exact decisions in
Section~\ref{sec:algorithms}, applied to representatives of these
universes, give the directly computed columns of
Table~\ref{tab:card}: connected graphs, Euler characteristic one,
$\mathbb F_2$-acyclicity, $\IS$, $\ISV$, the four dismantling levels
$0$ through $3$, and cones. The columns for $\I$ and integral
acyclicity are deduced below by classifying the remaining acyclic
graphs.

The finite enumerations and the resulting cardinalities in this
section are computer-assisted. The accompanying material records
the input universes, comparisons and finite verifications, with a map
from the results here to their supporting files.
The mathematical procedures establish what is counted; the retained
execution records supply their numerical results. The finite enumerations
use the generation and isomorphism software of McKay and Piperno
\cite{McKayPiperno2014}; the versions and generation restrictions are
recorded with the data.

\subsection{The first acyclic exceptions}

The directly computed acyclic and strong columns agree through order
$12$. At order $13$ their difference is three. 

\begin{theorem}\label{thm:classification13}
Up to isomorphism, the only $\mathbb F_2$-acyclic graphs on at most
$13$ vertices outside $\IS$ are $\Gt,\Gt',\Gt''$. Every graph in
$\I$ on at most $13$ vertices belongs to $\IS$, and $\Gt$ is the
only stuck graph of order $13$.
\end{theorem}
\begin{proof}
Every $\mathbb F_2$-acyclic graph is connected, so the connected census
contains both classes. Every strong graph is $\mathbb F_2$-acyclic, so equality of the two
finite cardinalities through order $12$ gives equality of the classes
in those orders. The cases of orders one and two are immediate.
At order $13$ the difference is
\[
 \RowThirteenAcyclic-\RowThirteenIS=3.
\]
Theorem~\ref{thm:A} supplies three acyclic non-strong graphs with
different edge counts, hence pairwise non-isomorphic. They exhaust
the difference. Their nonzero integral homology excludes them from
$\I$, which lies between $\IS$ and $\mathcal A_{\mathbb F_2}$.
Finally, the link calculations in Table~\ref{tab:witness-links}
show that $\Gt$ is stuck, whereas $\Gt'$ and $\Gt''$ each admit the
edge deletion reversing the last specified addition. Every stuck graph
is acyclic and non-strong, so there are no others at this order.
\end{proof}

\subsection{The \texorpdfstring{$K_4$}{K4}-free graphs of order fourteen}

The following classification supplies the positive exceptions in
the last row of the table. Its reduction uses only the complete
classification through order $13$ and the local $K_4$-free criterion.

\begin{theorem}\label{thm:classification14}
There are $51$ isomorphism classes of stuck $K_4$-free graphs of
order $14$. Exactly three lie in $\I$, namely $\Ga,\Gb,\Gc$;
each of the other $48$ has first integral homology $\ZZ/3$ and
zero higher homology. There are exactly eight $K_4$-free graphs
of order $14$ in $\I\setminus\IS$, up to isomorphism. Their
minimum number of edges is $44$, attained by $\Ga$.
\end{theorem}
\begin{proof}
First, enumerate the connected $K_4$-free graphs of orders two through
seven and apply Proposition~\ref{rem:K4-local}. No graph is almost
stuck. This check examines, respectively, $1,2,5,17,82,536$
isomorphism classes; it includes graphs with no admissible vertex.
If a stuck $K_4$-free graph of order at most $14$ had a cut vertex,
Lemma~\ref{lem:2conn} would give two almost stuck pieces with
orders summing to at most $15$. One would have order at most seven,
a contradiction. The seeds are therefore $2$-connected and satisfy
$4\le\delta\le\Delta\le12$, by Lemma~\ref{lem:degrees}.

Enumerating the $K_4$-free graphs of order $14$ with these restrictions,
and testing $\mathbb F_2$-homology and the local neighbourhood conditions, gives
$51$ pairwise non-isomorphic stuck graphs. Theorem~\ref{thm:B}
and Proposition~\ref{prop:other-seeds} prove that $\Ga,\Gb,\Gc$
lie in $\I$. Integral homology of the
other $48$, computed by~\eqref{eq:integral-homology}, has free
rank zero in degree one, and $\partial_2$ has a single nonunit Smith
factor, equal to $3$. Thus $H_1\cong\ZZ/3$. Their clique complexes
are connected, have dimension at most two and Euler characteristic
one. Hence $\widetilde H_0=0$, and $H_2$, which is free, has rank
zero. All higher homology therefore vanishes. In particular,
none of these $48$ seeds lies in $\I$.
The list of seeds, together with its homology
and the finite checks, is included in the accompanying repository.

By Theorem~\ref{thm:classification13}, a graph in
$\I\setminus\IS$ has at least $14$ vertices. Apply
Lemma~\ref{lem:stuck} to the $K_4$-free class: every $K_4$-free
member of $\I\setminus\IS$ of order $14$ reduces by edge deletions
to a stuck graph in $\I$.
It therefore arises from one of $\Ga,\Gb,\Gc$. Their pruned
edge-addition closures, restricted to $K_4$-free graphs, have the
following numbers of non-strong classes, grouped by edge count:
\[
\begin{array}{c|rrrr|r}
 \text{seed} & 44 & 45 & 46 & 47 & \text{total}\\ \hline
 \Ga & 1 & 2 & 1 & 0 & 4\\
 \Gb & 0 & 1 & 1 & 1 & 3\\
 \Gc & 0 & 1 & 0 & 0 & 1
\end{array}
\]
The three sets are pairwise disjoint up to isomorphism, so their
union has eight classes. Proposition~\ref{prop:closure-correct}
shows that no non-strong target is lost. Every retained graph belongs
to $\I$, since the additions are I-transformations; every possible
target has now been included. This proves the count of eight.

The three contractible seeds have $44,45,45$ edges. Edge additions
cannot decrease these numbers, so every target has at least $44$
edges. The graph $\Ga$ attains the bound.
\end{proof}

The eight graphs have one-edge extensions in $\IS$, each obtained by
an admissible addition. Thus their membership in $\I$ can also be
checked individually by~\eqref{eq:I-path}; the full lists and
reductions are supplied with the repository.

\subsection{Adjudication of the integral and contractible columns}

The coefficient distinction is essential to the first conjecture.
Within each finite universe, the inclusions
\begin{equation}\label{eq:table-chain}
 \IS\subseteq\I\subseteq\mathcal A_{\ZZ}
       \subseteq\mathcal A_{\mathbb F_2}
\end{equation}
leave only the acyclic non-strong graphs to be decided. Through
order $12$ there are none. At order $13$ all three have integral
torsion, as established above.

At order $14$, let $G$ be $K_4$-free, $\mathbb F_2$-acyclic and
outside $\IS$. An admissible deletion preserves its homotopy type
and cannot leave a strong graph, since that would give a strong
reduction of $G$. Delete admissible edges until either a vertex
deletion is possible or no admissible element remains. In the first
case the vertex deletion gives an acyclic non-strong graph of order
$13$. Theorem~\ref{thm:classification13} identifies it with one of
the three torsion graphs. Hence $G$ also has nonzero integral homology.
In the second case the terminal graph is one of the $51$ stuck seeds.
A torsion seed again gives nonzero integral homology. A contractible
seed gives one of the eight graphs of
Theorem~\ref{thm:classification14}. This exhausts the possibilities.
In every torsion case the first integral homology is $\ZZ/3$ and
the reduced homology vanishes in all other degrees.

\begin{proposition}\label{prop:table-integral}
In each universe of Table~\ref{tab:card}, integral acyclicity and
I-contractibility coincide. Their common cardinality equals the
$\mathbb F_2$-acyclic count through order $12$, that count minus
three at order $13$, and that count minus $625$ in the $K_4$-free
row of order $14$.
\end{proposition}
\begin{proof}
The preceding classification partitions every acyclic non-strong
graph into a member of $\I$ or a graph with nonzero integral
homology. Together with~\eqref{eq:table-chain}, this proves the
equality of classes in the specified universes. At order $14$,
the acyclic and strong counts differ by
\[
 \RowFourteenAcyclic-\RowFourteenIS=633.
\]
Exactly eight of these graphs lie in $\I$, leaving $625$ with
nonzero integral homology. The preceding reduction also determines their homology:
$\widetilde H_1\cong\ZZ/3$, with vanishing reduced homology in
all other degrees. Separate integral calculations on the $625$
graphs confirm these groups. Thus the integral entries are
$\RowThirteenI$ at order $13$ and $\RowFourteenI$ in the last
row. All earlier entries follow from equality of the acyclic and
strong classes.
\end{proof}

The equality just proved concerns the listed finite universes; it
does not identify integral acyclicity and contractibility in general.

Table~\ref{tab:card} combines these deductions with the direct counts.
The $63$ entries shared with the published table
\cite[Table~1]{DEFH2023} agree: these are the connected,
$\mathcal A_{\mathbb F_2}$, $\I$, $\IS$, $\ISV$, $1$-Dis
and $0$-Dis columns at orders $3$ through $11$.

\begin{table}[tbp]
\centering
\caption{Numbers of isomorphism classes of connected graphs.
The two panels give twelve counts, with $n$ repeated.
The entire row $14^{\dagger}$ is restricted to $K_4$-free graphs.
$\mathcal A_R$ denotes reduced acyclicity over $R$, and $(-1)$-Dis
denotes the cones.}\label{tab:card}
% Thirteen logical columns, including n; twelve counts in two panels.
\begingroup
\footnotesize
\setlength{\tabcolsep}{2pt}
\sisetup{group-separator={,}}
\renewcommand{\arraystretch}{1.12}
\begin{tabular*}{\textwidth}{@{\extracolsep{\fill}}rrrrrrr@{}}
\toprule
$n$ & connected & $\chi=1$ & $\mathcal A_{\mathbb F_2}$ & $\mathcal A_{\ZZ}$ & $\I$ & $\IS$ \\
\midrule
3 & \num{2} & \num{2} & \num{2} & \num{2} & \num{2} & \num{2} \\
4 & \num{6} & \num{5} & \num{5} & \num{5} & \num{5} & \num{5} \\
5 & \num{21} & \num{16} & \num{16} & \num{16} & \num{16} & \num{16} \\
6 & \num{112} & \num{68} & \num{68} & \num{68} & \num{68} & \num{68} \\
7 & \num{853} & \num{404} & \num{403} & \num{403} & \num{403} & \num{403} \\
8 & \num{11117} & \num{3837} & \num{3798} & \num{3798} & \num{3798} & \num{3798} \\
9 & \num{261080} & \num{68458} & \num{66167} & \num{66167} & \num{66167} & \num{66167} \\
10 & \num{11716571} & \num{2571263} & \num{2343148} & \num{2343148} & \num{2343148} & \num{2343148} \\
11 & \num{1006700565} & \num{204131202} & \num{166978401} & \num{166978401} & \num{166978401} & \num{166978401} \\
12 & \num{164059830476} & \num{32482574568} & \num{22666537431} & \num{22666537431} & \num{22666537431} & \num{22666537431} \\
13 & \num{50335907869219} & \num{9859955384222} & \num{5619022190286} & \num{5619022190283} & \num{5619022190283} & \num{5619022190283} \\
14$^{\dagger}$ & \num{309317860269934} & \num{6112638767338} & \num{755877673752} & \num{755877673127} & \num{755877673127} & \num{755877673119} \\
\bottomrule
\end{tabular*}
\par\medskip
\begin{tabular*}{\textwidth}{@{\extracolsep{\fill}}rrrrrrr@{}}
\toprule
$n$ & $\ISV$ & $3$-Dis & $2$-Dis & $1$-Dis & $0$-Dis & $(-1)$-Dis \\
\midrule
3 & \num{2} & \num{2} & \num{2} & \num{2} & \num{2} & \num{2} \\
4 & \num{5} & \num{5} & \num{5} & \num{5} & \num{5} & \num{4} \\
5 & \num{16} & \num{16} & \num{16} & \num{16} & \num{16} & \num{11} \\
6 & \num{68} & \num{68} & \num{68} & \num{68} & \num{68} & \num{34} \\
7 & \num{403} & \num{403} & \num{403} & \num{403} & \num{403} & \num{156} \\
8 & \num{3798} & \num{3798} & \num{3798} & \num{3798} & \num{3791} & \num{1044} \\
9 & \num{66167} & \num{66167} & \num{66167} & \num{66167} & \num{65561} & \num{12346} \\
10 & \num{2343148} & \num{2343148} & \num{2343148} & \num{2343148} & \num{2258314} & \num{274668} \\
11 & \num{166978389} & \num{166978389} & \num{166978389} & \num{166978389} & \num{151084672} & \num{12005168} \\
12 & \num{22666525713} & \num{22666525713} & \num{22666525713} & \num{22666525693} & \num{18510237804} & \num{1018997864} \\
13 & \num{5619007241580} & \num{5619007241580} & \num{5619007241580} & \num{5619007205823} & \num{4015494549880} & \num{165091172592} \\
14$^{\dagger}$ & \num{755266302709} & \num{755266302709} & \num{755266302709} & \num{755266302709} & \num{81156275851} & \num{20797002} \\
\bottomrule
\end{tabular*}
\endgroup

\end{table}

\subsection{Consequences of the table}

The constructions and the complete lower-order counts now determine
the three minima relevant to the conjectures.

\begin{corollary}[Computer-assisted]\label{cor:minimality}
The minimum order of an $\mathbb F_2$-acyclic graph outside $\IS$
is $13$. The minimum order of a graph in $\I\setminus\IS$ is
$14$, and the same minimum holds for an integrally acyclic graph
outside $\IS$. The minimum order of a strong graph having an
$\IS$-admissible edge whose deletion leaves a non-strong graph
is also $14$. The minimum number of vertices of a contractible,
non-collapsible flag simplicial complex is $14$.
\end{corollary}
\begin{proof}
Theorem~\ref{thm:classification13} excludes the smaller orders,
and Theorems~\ref{thm:A} and~\ref{thm:B} provide the first two
examples. Proposition~\ref{prop:table-integral} excludes an
integrally acyclic non-strong graph through order $13$, while
$\Ga\in\I$ attains order $14$. Finally, an admissible edge
deletion from a strong graph is an I-transformation. If its result
is non-strong, that result lies in $\I\setminus\IS$ and has the
same order. The lower bound is therefore $14$; the one-edge
extension of $\Ga$ in Section~\ref{sec:examples} attains it.

If a contractible flag complex $K=\cl(G)$ has at most $13$ vertices,
then $G\in\I$ by \cite[Corollary~1.5]{DochtermannMatsushita2025}.
Theorem~\ref{thm:classification13} gives $G\in\IS$, and
\cite[Theorem~1]{EFH2022} shows that $K$ is collapsible.
Proposition~\ref{prop:noncollapsible} supplies the example
$\cl(\Ga)$ on $14$ vertices.
\end{proof}

The dismantling columns give further separations. The classes
$2$-Dis and $3$-Dis have the same counts as $\ISV$ throughout the
table. The first difference between $2$-Dis and $1$-Dis occurs at
order $12$, where it has cardinality $20$; at order $13$ the
difference is $\NumTwoDisNotOneDisThirteen$. These conclusions
follow by subtracting the counts of nested classes. The $20$ graphs
of order $12$ are catalogued in the accompanying material. The first
difference between $\IS$ and $\ISV$, already established in
\cite[Theorem~11]{DEFH2023}, occurs at order $11$ and consists of
twelve graphs.

There is also a direct interpretation of the cone column. Adjoining
a universal vertex is a bijection between isomorphism classes of
graphs on $n-1$ vertices and cones on $n$ vertices. Its inverse is
well defined because deleting any of several universal vertices
gives isomorphic graphs. This explains the agreement of the
unrestricted cone entries with the known numbers of unlabelled
graphs~\cite{OEISA000088}; the connected totals at orders $3$ through
$13$ agree with~\cite{OEISA001349}. A cone is $K_4$-free exactly when deletion
of its apex leaves a triangle-free graph. The last cone entry is
therefore $\RowFourteenCone$, the number of triangle-free graphs
on $13$ vertices~\cite{OEISA006785}.

\appendix
\section{Incidence data and reductions of the witnesses}\label{app:certificates}

The tables below specify the four graphs used in the constructions
and the finite calculations in their proofs. In particular, their
definitions do not depend on recovering edges from a drawing.

\begin{table}[htbp]
\centering\small
\caption{For each vertex $i$, the entries are precisely its neighbours
larger than $i$. Thus each edge occurs once. The symbol $\varnothing$
means that there is no larger neighbour; the dash indicates that the
vertex is absent.}\label{tab:witness-adjacency}
\begin{tabular}{@{}rllll@{}}
\toprule
$i$ & $\Gt$ & $\Ga$ & $\Gb$ & $\Gc$ \\
\midrule
0 & $4,5,9,10,11,12$ & $4,5,8,10,12,13$ & $5,7,9,12,13$ & $4,5,7,9,12,13$ \\
1 & $4,6,8,9,11,12$ & $5,7,8,10,12,13$ & $5,8,9,11,13$ & $4,5,8,10,12,13$ \\
2 & $5,7,8,9,10,12$ & $6,7,9,12,13$ & $6,7,8,11,12,13$ & $6,7,8,10,11,13$ \\
3 & $6,7,8,10,11,12$ & $6,9,10,12,13$ & $6,7,9,10,12,13$ & $6,7,9,10,11,13$ \\
4 & $8,10,12$ & $8,9,11,13$ & $6,8,9,10,11,13$ & $8,9,11,13$ \\
5 & $8,11,12$ & $9,10,11,12$ & $7,8,10,11,12$ & $9,10,11,13$ \\
6 & $9,10,12$ & $7,10,11,12$ & $10,11,12$ & $8,9,12,13$ \\
7 & $9,11,12$ & $8,11,13$ & $10,11,13$ & $10,12,13$ \\
8 & $10,11$ & $11,12$ & $10,12,13$ & $11,12$ \\
9 & $10,11$ & $11,12,13$ & $11,12,13$ & $11,12$ \\
10 & $\varnothing$ & $11,13$ & $\varnothing$ & $11,12$ \\
11 & $\varnothing$ & $\varnothing$ & $\varnothing$ & $\varnothing$ \\
12 & $\varnothing$ & $\varnothing$ & $\varnothing$ & $\varnothing$ \\
13 & --- & $\varnothing$ & $\varnothing$ & $\varnothing$ \\
\bottomrule
\end{tabular}

\end{table}

\subsection{A triangulated disc for the thirteen-vertex graph}
\label{app:disc-certificates}

Here is an incidence description of the disc in Figure~\ref{fig:G13}.
Its vertices are $0,1,\ldots,7,12$ and $b_0,\ldots,b_{11}$; its
triangles are listed in Table~\ref{tab:G13-disc-triangles}.
In the figure, $b_0$ is the topmost boundary vertex, and the boundary
indices increase clockwise.
The interior vertices keep their labels. On the boundary, the map
to $\Gt$ is
\[
 b_{4r}\longmapsto10,\qquad
 b_{4r+1}\longmapsto9,\qquad
 b_{4r+2}\longmapsto11,\qquad
 b_{4r+3}\longmapsto8\qquad(0\le r\le2).
\]

\begin{table}[htbp]
\centering\small
\caption{The triangles of the disc whose boundary is attached with
degree three.}\label{tab:G13-disc-triangles}
\begin{tabular}{@{}rlrlrlrl@{}}
\toprule
$j$ & $T_j$ & $j$ & $T_j$ & $j$ & $T_j$ & $j$ & $T_j$ \\
\midrule
1 & $\{0,4,b_{4}\}$ & 8 & $\{1,4,12\}$ & 15 & $\{2,7,b_{9}\}$ & 22 & $\{3,7,12\}$ \\
2 & $\{0,4,12\}$ & 9 & $\{1,6,b_{1}\}$ & 16 & $\{2,7,12\}$ & 23 & $\{3,b_{11},b_{0}\}$ \\
3 & $\{0,5,b_{6}\}$ & 10 & $\{1,6,12\}$ & 17 & $\{2,b_{7},b_{8}\}$ & 24 & $\{3,b_{11},b_{10}\}$ \\
4 & $\{0,5,12\}$ & 11 & $\{1,b_{3},b_{2}\}$ & 18 & $\{2,b_{9},b_{8}\}$ & 25 & $\{4,b_{3},b_{4}\}$ \\
5 & $\{0,b_{5},b_{4}\}$ & 12 & $\{1,b_{1},b_{2}\}$ & 19 & $\{3,6,b_{0}\}$ & 26 & $\{5,b_{7},b_{6}\}$ \\
6 & $\{0,b_{5},b_{6}\}$ & 13 & $\{2,5,b_{7}\}$ & 20 & $\{3,6,12\}$ & 27 & $\{6,b_{1},b_{0}\}$ \\
7 & $\{1,4,b_{3}\}$ & 14 & $\{2,5,12\}$ & 21 & $\{3,7,b_{10}\}$ & 28 & $\{7,b_{9},b_{10}\}$ \\
\bottomrule
\end{tabular}

\end{table}

To verify the disc assertion, start with $T_1$ and attach the remaining
triangles in this order:
\[
\begin{gathered}
2,4,3,5,6,8,7,10,9,11,12,14,13,16,15,17,18,\\
20,19,22,21,23,24,25,26,27,28.
\end{gathered}
\]
At each step, the intersection with the preceding union is either
one boundary edge or two consecutive boundary edges, with no
additional vertices. This follows directly from the triangle list.
Attaching a triangle along either such boundary arc preserves a
closed disc. The final boundary is
$b_0,b_1,\ldots,b_{11},b_0$.

The triangles map bijectively onto the $28$ triangles of $\cl(\Gt)$.
Their edges give exactly the adjacency data for $\Gt$, which contain
no $K_4$, so the image is the entire clique complex. Each interior
vertex has a distinct label outside $\{8,9,10,11\}$. The images of
interior edges are distinct, and each has an endpoint outside the
four-cycle. Thus only equal-labelled boundary vertices and their
corresponding boundary edges are identified. This verifies the
hypotheses of Proposition~\ref{prop:disc}, with boundary degree three.

\subsection{The obstruction to a first deletion}

The graphs in Table~\ref{tab:witness-links} are $K_4$-free and have
connected vertex neighbourhoods. The table gives their first Betti
numbers $|E(\Nb_G(v))|-|V(\Nb_G(v))|+1$ in vertex-label order.
All are positive; every edge neighbourhood is independent with at
least two vertices. Hence all these neighbourhoods have nonzero
reduced homology and none belongs to $\IS$.

\begin{table}[htbp]
\centering\small
\caption{Nonzero homology of the vertex and edge neighbourhoods.
The Betti numbers of vertex neighbourhoods have the same values over
every field.}\label{tab:witness-links}
\begin{tabular}{@{}lll@{}}
\toprule
Graph & First Betti numbers of vertex neighbourhoods & Edge-neighbourhood sizes\\
\midrule
$\Gt$ & $1,1,1,1,1,1,1,1,2,2,2,2,1$ & $2,3$\\
$\Ga$ & $1,1,1,1,1,2,1,1,1,2,2,1,2,2$ & $2,3$\\
$\Gb$ & $1,1,1,1,1,2,1,2,2,2,1,1,1,3$ & $2,3,4$\\
$\Gc$ & $2,1,1,1,1,1,1,2,1,2,2,2,1,2$ & $2,3$\\
\bottomrule
\end{tabular}
\end{table}

\subsection{Reductions of the one-edge extensions}

Tables~\ref{tab:Ga-reduction} and~\ref{tab:variant-reductions} give
the local neighbourhood at every deletion. A displayed path specifies
exactly its consecutive edges. A single label denotes a one-vertex
graph. At step $13$ for $\Gb+\{0,3\}$, the two displayed paths are
united at their common vertices; their union is a tree.
All neighbourhoods are evaluated before the corresponding deletion,
and all three reductions end at the vertex $13$.

\begin{table}[htbp]
\centering\small
\caption{The complete reduction of $\Ga+\{0,1\}$ used in
Theorem~\ref{thm:B}.}\label{tab:Ga-reduction}
\begin{tabular}{@{}rll@{}}
\toprule
Step & Delete & Current neighbourhood \\
\midrule
1 & $e(0,5)$ & $10\text{--}1\text{--}12$ \\
2 & $e(0,10)$ & $1\text{--}13$ \\
3 & $e(0,12)$ & $1\text{--}8$ \\
4 & $e(8,12)$ & $1$ \\
5 & $e(1,12)$ & $5$ \\
6 & $v(5)$ & $1\text{--}10\text{--}11\text{--}9\text{--}12$ \\
7 & $v(10)$ & $1\text{--}13\text{--}3\text{--}6\text{--}11$ \\
8 & $v(3)$ & $6\text{--}12\text{--}9\text{--}13$ \\
9 & $v(6)$ & $11\text{--}7\text{--}2\text{--}12$ \\
10 & $v(2)$ & $7\text{--}13\text{--}9\text{--}12$ \\
11 & $v(7)$ & $11\text{--}8\text{--}1\text{--}13$ \\
12 & $v(1)$ & $8\text{--}0\text{--}13$ \\
13 & $v(0)$ & $8\text{--}4\text{--}13$ \\
14 & $v(4)$ & $8\text{--}11\text{--}9\text{--}13$ \\
15 & $v(8)$ & $11$ \\
16 & $v(11)$ & $9$ \\
17 & $v(12)$ & $9$ \\
18 & $v(9)$ & $13$ \\
\bottomrule
\end{tabular}

\end{table}

\begin{table}[htbp]
\centering\small
\caption{Reductions of the other two contractible seeds after one
admissible edge addition. A dash means that the singleton has already
been reached.}\label{tab:variant-reductions}
\begin{tabular}{@{}rllll@{}}
\toprule
& \multicolumn{2}{c}{$\Gb+\{0,3\}$} & \multicolumn{2}{c}{$\Gc+\{2,3\}$} \\
Step & Delete & Current neighbourhood & Delete & Current neighbourhood \\
\midrule
1 & $e(0,9)$ & $12\text{--}3\text{--}13$ & $e(2,7)$ & $10\text{--}3\text{--}13$ \\
2 & $e(0,13)$ & $3\text{--}7$ & $e(2,10)$ & $3\text{--}11$ \\
3 & $e(9,12)$ & $3$ & $e(2,13)$ & $3\text{--}6$ \\
4 & $e(3,9)$ & $13$ & $e(6,13)$ & $3$ \\
5 & $e(3,13)$ & $7$ & $e(3,13)$ & $7$ \\
6 & $e(7,13)$ & $2$ & $v(7)$ & $3\text{--}10\text{--}12\text{--}0\text{--}13$ \\
7 & $v(2)$ & $7\text{--}11\text{--}6\text{--}12\text{--}8\text{--}13$ & $v(10)$ & $3\text{--}11\text{--}5\text{--}1\text{--}12$ \\
8 & $v(6)$ & $11\text{--}4\text{--}10\text{--}3\text{--}12$ & $v(1)$ & $5\text{--}13\text{--}4\text{--}8\text{--}12$ \\
9 & $v(3)$ & $10\text{--}7\text{--}0\text{--}12$ & $v(4)$ & $8\text{--}11\text{--}9\text{--}0\text{--}13$ \\
10 & $v(0)$ & $7\text{--}5\text{--}12$ & $v(0)$ & $12\text{--}9\text{--}5\text{--}13$ \\
11 & $v(4)$ & $10\text{--}8\text{--}13\text{--}9\text{--}11$ & $v(8)$ & $11\text{--}2\text{--}6\text{--}12$ \\
12 & $v(7)$ & $10\text{--}5\text{--}11$ & $v(2)$ & $6\text{--}3\text{--}11$ \\
13 & $v(5)$ & $11\text{--}1\text{--}8\text{--}10,\ 8\text{--}12$ & $v(3)$ & $6\text{--}9\text{--}11$ \\
14 & $v(1)$ & $8\text{--}13\text{--}9\text{--}11$ & $v(6)$ & $9\text{--}12$ \\
15 & $v(10)$ & $8$ & $v(11)$ & $5\text{--}9$ \\
16 & $v(11)$ & $9$ & $v(12)$ & $9$ \\
17 & $v(9)$ & $13$ & $v(9)$ & $5$ \\
18 & $v(12)$ & $8$ & $v(5)$ & $13$ \\
19 & $v(8)$ & $13$ & --- & --- \\
\bottomrule
\end{tabular}

\end{table}

\FloatBarrier
\section*{Computational material}
The accompanying reproducibility package contains the implementations,
catalogues of exceptional graphs, reduction sequences, homology computations
and enumeration records supporting the finite results. It provides
local checks of the examples and instructions for reproducing the
complete censuses. The project repository is
\href{https://github.com/Geometric-and-Combinatorial-Structures/minimal-I-contractible-not-strong}{\texttt{minimal-I-contractible-not-strong}}.

\section*{Acknowledgements}
The authors thank the \'Area de C\'omputo de Alto Rendimiento de la
Universidad de Sonora (ACARUS) for computing time on the Yuca cluster
and for technical support.

\section*{Declarations}
\noindent\textbf{Funding.}
This work was supported by the Universidad de Sonora through the internal
research funding call of the Direcci\'on de Apoyo a Docentes,
Investigaci\'on y Posgrado (DADIP), project USO315010378.

\noindent\textbf{Author contributions.}
The authors are listed alphabetically and contributed equally to this work.

\noindent\textbf{Competing interests.}
The authors declare that they have no known competing financial interests
or personal relationships that could have appeared to influence the work
reported in this paper.

\noindent\textbf{Use of AI-assisted tools.}
Codex (OpenAI) and Claude (Anthropic) were used to draft and revise
mathematical exposition, arguments, tables and software, and to assist
with source checking and adversarial review. The authors retain
responsibility for the mathematical content and the final manuscript.

\bibliographystyle{siam}
\bibliography{refs}
\end{document}
